
\chapter{Dynamic Logic of Relative Responses and Trajectories}
\label{ch:dynamic-relative-response-trajectories}

\section{Orientation: spans, profiles, and finite trajectories}
\label{sec:dynamic-response-trajectory-orientation}

The preceding two chapters constructed the same relative Mealy semantics in
three compatible forms.  A Mealy morphism
\[
    \Phi=(P,\delta_P,\omega_P):
    \mathbb A\rightsquigarrow\mathbb B
\]
and a right \(\mathbb B\)-action
\[
    Y=(Y\xrightarrow{r}B_0,\sigma)
\]
determine the relative state object
\[
    S_{\Phi,Y}=P\times_{B_0}Y,
\]
the relative execution category
\[
    \mathsf{Prog}_{\Phi,Y},
\]
the multiplicative response coalgebra
\[
    c_{\Phi,Y}:
    S_{\Phi,Y}\to
    \mathcal M_{\mathbb A,\mathbb B,Y}(S_{\Phi,Y}),
\]
and the finite trajectory object
\[
    \mathscr T_{\Phi,Y}
    :=
    \operatorname{Path}(\mathsf{Prog}_{\Phi,Y}).
\]
The present chapter develops the corresponding logical semantics.

There are three logically distinct levels.
\[
\begin{array}{c|c|c}
\text{level} & \text{semantic object} & \text{quantification}\\
\hline
\text{span} &
X\xleftarrow{\ell}R\xrightarrow{r}Y &
\text{selected execution witnesses}\\
\text{profile} &
c:X\to\mathcal M(X) &
\text{positions in a total response profile}\\
\text{trajectory} &
\mathscr T(n) &
\text{finite segmented executions}
\end{array}
\]
At the span level, diamonds and boxes quantify over a chosen execution span.
At the profile level, predicate liftings inspect the positions of a total
response policy before the coalgebra is flattened.  At the trajectory level,
the same predicate-transformer calculus is applied to the initial and final
state maps of finite path objects, or to any of their intermediate boundary
maps.

The chapter proves two exact comparison results.  First, guarded existential
profile liftings pull back along \(c_{\Phi,Y}\) to the diamonds of the
corresponding guarded one-step spans.  Second, weak universal profile boxes
pull back to the boxes of those restricted spans.  The strong profile box is
different: it asserts total guarded responsiveness for every input position,
rather than merely quantifying over witnesses already selected by a
restriction.

The finite-trajectory layer adds a different distinction.  Endpoint-only
modalities for unrestricted categorical paths contract to one-step
modalities, because a single arrow of \(\mathbb A\) may already be composite.
Boundary-sensitive formulas and generator-restricted path formulas do not
collapse in this way.  They retain information about the chosen segmentation
of an execution.

\subsection{Standing hypotheses and levels of logical structure}
\label{subsec:dynamic-logical-hypotheses}

The constructions are stratified by categorical strength.
\[
\begin{array}{c|c}
\text{construction} & \text{assumption}\\
\hline
\text{pullback of predicates and tests} & \text{finite limits}\\
\text{diamonds} & \text{regularity}\\
\text{finite guarded joins} & \text{coherence}\\
\text{boxes and intuitionistic implication} &
\text{first-order Heyting structure}\\
\text{all-horizon safety meets} &
\text{the displayed small meets in subobject lattices}
\end{array}
\]
No dependent products used to construct the response polynomial are repeated
here.  They enter only through the already constructed coalgebra
\(c_{\Phi,Y}\) and its position object.  Whenever a theorem uses
Beck--Chevalley for existential or universal quantification, that hypothesis
is stated explicitly.

\subsection{Notation for relative execution}
\label{subsec:dynamic-relative-execution-notation}

Write
\[
    S:=S_{\Phi,Y},
    \qquad
    M:=M_{\Phi,Y}.
\]
The one-step execution span is
\[
    \mathbf m_{\Phi,Y}
    =
    \left(
    S\xleftarrow{s_{\Phi,Y}}M
    \xrightarrow{t_{\Phi,Y}}S
    \right).
\]
Its canonical observations are
\[
    \iota_A:M\to A_1,
    \qquad
    \iota_B:M\to B_1,
    \qquad
    \lambda_{\Phi,Y}:
    M\to\left(\int_{\mathbb B}Y\right)_1.
\]
They record, respectively, the input arrow, the raw output arrow, and the
downstream transition.  The response coalgebra and its universal position
object are those of
Chapter~\ref{ch:relative-mealy-evaluation-polynomial-response}.  The finite
trajectory objects and their Segal operations are those of
Chapter~\ref{ch:temporal-trajectories-segal-execution}.

The sequential-composition notation is the dynamic-logic notation:
\[
    R;S
\]
means that \(R\) is executed first and \(S\) second.

\section{Predicate transformers for spans}
\label{sec:relative-predicate-transformers}

We now pass from the operational layer to the logical layer. The notation is
chosen to match the non-iterative dynamic-logic reading of programs as
predicate transformers: diamonds express existence of a successful execution,
boxes express universal correctness over all executions, tests restrict
execution to states satisfying a predicate, and finite choice denotes
nondeterministic branching \cite{HarelKozenTiurynDynamicLogic}. The
categorical implementation uses subobjects as predicates, regular images as
existential quantification, and right adjoints to pullback as universal
quantification \cite{Jacobs1999CategoricalLogic}. In the regular fragment,
this agrees with the usual relational reading of programs, where relations can
be represented categorically by suitable spans
\cite{FongSpivakRegularRelational}.

\subsection{Regular images and existential quantification}
\label{subsec:relative-regular-images}

\begin{definition}[Regular category]
\label{def:relative-regular-category}
A category \(\mathcal E\) is \textbf{regular} if it has finite limits, every
morphism factors as a regular epimorphism followed by a monomorphism, and
regular epimorphisms are stable under pullback.
\end{definition}

For the rest of this section, assume \(\mathcal E\) is regular. For an object
\(X\), write
\[
\mathrm{Sub}(X)
\]
for the poset of subobjects of \(X\). For a morphism \(f:X\to Y\), write
\[
f^\ast:\mathrm{Sub}(Y)\to\mathrm{Sub}(X)
\]
for pullback.

\begin{definition}[Existential image]
\label{def:relative-existential-image}
For a morphism
\[
f:X\to Y,
\]
define
\[
\exists_f:\mathrm{Sub}(X)\to\mathrm{Sub}(Y)
\]
by sending a subobject \(U\hookrightarrow X\) to the regular image of
\[
U\to X\xrightarrow{f}Y.
\]
\end{definition}

\begin{proposition}[Existential adjunction]
\label{prop:relative-existential-adjunction}
For every morphism \(f:X\to Y\) in a regular category,
\[
\exists_f\dashv f^\ast.
\]
Thus, for \(U\in\mathrm{Sub}(X)\) and \(V\in\mathrm{Sub}(Y)\),
\[
\exists_f(U)\le V
\quad\Longleftrightarrow\quad
U\le f^\ast V.
\]
\end{proposition}

\begin{proof}
\leavevmode
\begin{enumerate}
\item \textbf{Translate the left-hand inequality.}
The statement
\[
\exists_f(U)\le V
\]
says that the image of the composite
\[
U\to X\xrightarrow{f}Y
\]
factors through \(V\hookrightarrow Y\).

\item \textbf{Remove the image.}
This is equivalent to the composite \(U\to Y\) itself factoring through
\(V\hookrightarrow Y\).

\item \textbf{Use the pullback.}
Factoring through \(V\) is equivalent to \(U\to X\) factoring through
\(f^\ast V\hookrightarrow X\).
\end{enumerate}
\end{proof}

\begin{proposition}[Beck--Chevalley for regular images]
\label{prop:relative-regular-beck-chevalley}
For every pullback square
\[
\begin{tikzcd}
X' \arrow[r,"g'"] \arrow[d,"f'"']
&
X \arrow[d,"f"]
\\
Y' \arrow[r,"g"']
&
Y
\arrow[ul,phantom,very near start,"\lrcorner"]
\end{tikzcd}
\]
in a regular category,
\[
g^\ast\exists_f=\exists_{f'}(g')^\ast.
\]
\end{proposition}

\begin{proof}
\leavevmode
\begin{enumerate}
\item \textbf{Factor the relevant composite.}
The regular image of a subobject after applying \(f\) is represented by a
regular epimorphism followed by a monomorphism.

\item \textbf{Pull the factorization back.}
Regular epimorphisms are stable under pullback, so pulling this image
factorization back along \(g\) gives the image factorization of the pulled
back subobject along \(f'\).

\item \textbf{Compare the two subobjects.}
The two constructions therefore determine the same subobject of \(Y'\).
\end{enumerate}
\end{proof}

When
\[
R:X\nrightarrow Y
\qquad\text{and}\qquad
S:Y\nrightarrow Z
\]
are spans, we write
\[
R;S:X\nrightarrow Z
\]
for the composite span obtained by performing \(R\) first and \(S\) second.
Equivalently, in the notation of
Chapter~\ref{ch:spans-internal-categories-mealy},
\[
R;S=S\circ R.
\]
This convention is the dynamic-logic convention for sequential composition of
programs.

\subsection{Diamond modalities}
\label{subsec:relative-diamond-modalities}

\begin{definition}[Diamond modality of a span]
\label{def:relative-span-diamond}
Let
\[
R=
\left(
X\xleftarrow{\ell}R\xrightarrow{r}Y
\right)
\]
be a span in a regular category. Define
\[
\langle R\rangle:\mathrm{Sub}(Y)\to\mathrm{Sub}(X)
\]
by
\[
\langle R\rangle\varphi
:=
\exists_\ell(r^\ast\varphi).
\]
The construction is represented by
\[
\begin{tikzcd}
r^\ast\varphi
\arrow[r]
\arrow[d]
&
\varphi
\arrow[d,hook]
\\
R
\arrow[r,"r"']
\arrow[d,"\ell"]
&
Y
\\
X
\end{tikzcd}
.
\]
\end{definition}

\begin{proposition}[Functoriality of diamonds]
\label{prop:relative-diamond-functoriality}
For spans
\[
R:X\nrightarrow Y
\qquad\text{and}\qquad
S:Y\nrightarrow Z
\]
in a regular category,
\[
\langle 1_X\rangle=\mathrm{id}_{\mathrm{Sub}(X)}
\]
and
\[
\langle R;S\rangle
=
\langle R\rangle\circ\langle S\rangle.
\]
\end{proposition}

\begin{proof}
\leavevmode
\begin{enumerate}
\item \textbf{Check the identity span.}
For the identity span
\[
X\xleftarrow{1_X}X\xrightarrow{1_X}X,
\]
the formula gives
\[
\exists_{1_X}1_X^\ast=\mathrm{id}_{\mathrm{Sub}(X)}.
\]

\item \textbf{Display the composite span.}
Write
\[
R=(X\xleftarrow{\ell}R\xrightarrow{r}Y),
\qquad
S=(Y\xleftarrow{m}S\xrightarrow{n}Z).
\]
Their composite has apex \(R\times_Y S\):
\[
\begin{tikzcd}[column sep=large,row sep=large]
R\times_Y S
  \arrow[r,"\pi_S"]
  \arrow[d,"\pi_R"']
&
S
  \arrow[d,"m"]
  \arrow[dr,"n",bend left=12]
&
\\
R
  \arrow[r,"r"']
  \arrow[d,"\ell"']
&
Y
&
Z
\\
X
&&
\end{tikzcd}
.
\]

\item \textbf{Compute the diamond of the composite.}
For \(\varphi\in\mathrm{Sub}(Z)\),
\[
\langle R;S\rangle\varphi
=
\exists_{\ell\pi_R}\bigl((n\pi_S)^\ast\varphi\bigr).
\]

\item \textbf{Use Beck--Chevalley.}
By functoriality of existential images and Beck--Chevalley for the pullback
defining \(R\times_Y S\), this is
\[
\exists_\ell r^\ast\exists_m n^\ast\varphi
=
\langle R\rangle(\langle S\rangle\varphi).
\]
\end{enumerate}
\end{proof}


\subsection{First-order Heyting structure}
\label{subsec:relative-first-order-heyting}

\begin{definition}[First-order Heyting category]
\label{def:relative-first-order-heyting}
A \textbf{first-order Heyting category} is a coherent category
\(\mathcal E\) such that, for every morphism
\[
f:X\to Y,
\]
the pullback functor
\[
f^\ast:\mathrm{Sub}(Y)\to\mathrm{Sub}(X)
\]
has a right adjoint
\[
\forall_f:\mathrm{Sub}(X)\to\mathrm{Sub}(Y),
\]
and these right adjoints satisfy Beck--Chevalley for pullback squares.
\end{definition}

Assume \(\mathcal E\) is a first-order Heyting category.

\subsection{Box modalities}
\label{subsec:relative-box-modalities}

\begin{definition}[Box modality of a span]
\label{def:relative-span-box}
For a span
\[
R=
(X\xleftarrow{\ell}R\xrightarrow{r}Y),
\]
define
\[
[R]:\mathrm{Sub}(Y)\to\mathrm{Sub}(X)
\]
by
\[
[R]\varphi:=\forall_\ell(r^\ast\varphi).
\]
\end{definition}

\begin{proposition}[Functoriality of boxes]
\label{prop:relative-box-functoriality}
For spans
\[
R:X\nrightarrow Y
\qquad\text{and}\qquad
S:Y\nrightarrow Z,
\]
one has
\[
[1_X]=\mathrm{id}_{\mathrm{Sub}(X)}
\]
and
\[
[R;S]=[R]\circ[S].
\]
\end{proposition}

\begin{proof}
\leavevmode
\begin{enumerate}
\item \textbf{Check the identity span.}
For the identity span, the formula gives
\[
\forall_{1_X}1_X^\ast=\mathrm{id}_{\mathrm{Sub}(X)}.
\]

\item \textbf{Use the composite pullback.}
For
\[
R=(X\xleftarrow{\ell}R\xrightarrow{r}Y),
\qquad
S=(Y\xleftarrow{m}S\xrightarrow{n}Z),
\]
the composite has apex \(R\times_Y S\).

\item \textbf{Apply Beck--Chevalley for \(\forall\).}
The Beck--Chevalley law for the pullback square defining \(R\times_Y S\),
together with functoriality of right adjoints, gives
\[
[R;S]\varphi
=
\forall_\ell r^\ast\forall_m n^\ast\varphi
=
[R]([S]\varphi).
\]
\end{enumerate}
\end{proof}


\section{Tests and state restrictions}
\label{subsec:relative-tests}

\begin{definition}[Tests]
\label{def:relative-tests}
For a state predicate
\[
\theta:\Theta\hookrightarrow S,
\]
define the corresponding test span
\[
\theta?
=
(S\xleftarrow{\theta}\Theta\xrightarrow{\theta}S).
\]
It is displayed as
\[
\begin{tikzcd}
& \Theta \arrow[dl,"\theta"'] \arrow[dr,"\theta"] & \\
S && S
\end{tikzcd}
.
\]
\end{definition}

\begin{proposition}[Diamond law for tests]
\label{prop:relative-test-diamond}
For every predicate \(\varphi\hookrightarrow S\),
\[
\langle\theta?\rangle\varphi
=
\theta\wedge\varphi.
\]
\end{proposition}

\begin{proof}
\leavevmode
\begin{enumerate}
\item \textbf{Expand the diamond.}
By definition,
\[
\langle\theta?\rangle\varphi
=
\exists_\theta(\theta^\ast\varphi).
\]

\item \textbf{Use monicity of the test.}
Since \(\theta\) is monic, existential image along \(\theta\) identifies
subobjects of \(\Theta\) with subobjects of \(S\) lying below \(\theta\).

\item \textbf{Identify the image.}
The image of \(\theta^\ast\varphi\) is therefore the meet
\[
\theta\wedge\varphi.
\]
\end{enumerate}
\end{proof}


\subsection{Boxes for tests}
\label{subsec:relative-boxes-tests}

\begin{proposition}[Box law for tests]
\label{prop:relative-test-box}
For a test
\[
\theta?=(S\xleftarrow{\theta}\Theta\xrightarrow{\theta}S),
\]
one has
\[
[\theta?]\varphi
=
\theta\Rightarrow\varphi.
\]
\end{proposition}

\begin{proof}
\leavevmode
\begin{enumerate}
\item \textbf{Expand the box.}
By definition,
\[
[\theta?]\varphi
=
\forall_\theta(\theta^\ast\varphi).
\]

\item \textbf{Test against an arbitrary predicate.}
For any predicate \(\psi\hookrightarrow S\),
\[
\psi\le \forall_\theta(\theta^\ast\varphi)
\]
if and only if
\[
\theta^\ast\psi\le \theta^\ast\varphi.
\]

\item \textbf{Use monicity of \(\theta\).}
Since \(\theta\) is monic, this is equivalent to
\[
\psi\wedge\theta\le\varphi.
\]

\item \textbf{Use the Heyting implication.}
The last inequality is equivalent to
\[
\psi\le\theta\Rightarrow\varphi.
\]
Hence
\[
[\theta?]\varphi=\theta\Rightarrow\varphi.
\]
\end{enumerate}
\end{proof}


\section{Guard loci and guarded one-step programs}
\label{sec:poly-relative-guard-loci-primitives}

\subsection{Guard loci}
\label{subsec:poly-relative-guard-loci}

Fix \(\Phi\) and \(Y\), and write
\[
S:=S_{\Phi,Y},
\qquad
M:=M_{\Phi,Y}.
\]

A guarded one-step specification may involve the following predicates:
\[
\theta\hookrightarrow S
\qquad
\text{current coupled-state guard},
\]
\[
R\hookrightarrow A_1
\qquad
\text{input-arrow guard},
\]
\[
O\hookrightarrow B_1
\qquad
\text{raw output-arrow guard},
\]
\[
T\hookrightarrow\left(\int_{\mathbb B}Y\right)_1
\qquad
\text{downstream transition guard},
\]
and
\[
\varphi\hookrightarrow S
\qquad
\text{continuation-state predicate}.
\]

The one-step execution object has canonical maps
\[
\iota_A:M\to A_1,
\qquad
\iota_A(a,x,y)=a,
\]
\[
\iota_B:M\to B_1,
\qquad
\iota_B(a,x,y)=\omega_P(a,x),
\]
and
\[
\lambda_{\Phi,Y}:M\to\left(\int_{\mathbb B}Y\right)_1,
\qquad
\lambda_{\Phi,Y}(a,x,y)=(\omega_P(a,x),y).
\]
They are displayed by
\[
\begin{tikzcd}[column sep=large,row sep=large]
&
M
  \arrow[dl,"\iota_A"']
  \arrow[d,"\lambda_{\Phi,Y}" description]
  \arrow[dr,"\iota_B"]
&
\\
A_1
&
\left(\int_{\mathbb B}Y\right)_1
  \arrow[l,phantom,"\mapsto"]
  \arrow[r,"\pi_{B_1}"']
&
B_1 .
\end{tikzcd}
\]

In the terminal downstream case,
\[
\int_{\mathbb B}\mathbf 1_{\mathbb B}\cong\mathbb B^{\mathrm{op}},
\]
so downstream transition guards reduce to predicates on output arrows of
\(\mathbb B\).

\subsection{Guarded one-step programs}
\label{subsec:poly-relative-guarded-primitives}

Let
\[
\theta\hookrightarrow S,
\qquad
R\hookrightarrow A_1,
\qquad
O\hookrightarrow B_1,
\qquad
T\hookrightarrow\left(\int_{\mathbb B}Y\right)_1
\]
be guards.

Pull them back to \(M\):
\[
\theta^\sharp:=s_{\Phi,Y}^\ast\theta,
\]
\[
R^\sharp:=\iota_A^\ast R,
\qquad
O^\sharp:=\iota_B^\ast O,
\qquad
T^\sharp:=\lambda_{\Phi,Y}^\ast T.
\]
Their simultaneous restriction is
\[
(\theta;R/O/T)^\sharp
:=
\theta^\sharp\wedge R^\sharp\wedge O^\sharp\wedge T^\sharp
\hookrightarrow M.
\]

\begin{definition}[Guarded one-step relative program]
\label{def:poly-relative-guarded-primitive}
The \textbf{guarded one-step relative program}
\[
\theta;R/O/T
\]
is the restricted span
\[
S
\xleftarrow{s_{\Phi,Y}i}
(\theta;R/O/T)^\sharp
\xrightarrow{t_{\Phi,Y}i}
S,
\]
where
\[
i:(\theta;R/O/T)^\sharp\hookrightarrow M
\]
is the inclusion.  Its diamond is denoted
\[
\langle\theta;R/O/T\rangle\varphi.
\]
\end{definition}

The restricted span is displayed by
\[
\begin{tikzcd}[column sep=large,row sep=large]
&
(\theta;R/O/T)^\sharp
  \arrow[dl,"s_{\Phi,Y}i"']
  \arrow[dr,"t_{\Phi,Y}i"]
  \arrow[d,hook,"i" description]
&
\\
S
&
M
  \arrow[l,"s_{\Phi,Y}"]
  \arrow[r,"t_{\Phi,Y}"']
&
S .
\end{tikzcd}
\]

\begin{proposition}[External reading of guarded one-step programs]
\label{prop:poly-relative-guarded-primitive-external}
In \(\mathbf{Set}\),
\[
(x,y)\models\langle\theta;R/O/T\rangle\varphi
\]
holds precisely when
\[
\theta(x,y)
\]
and there exists an admissible input arrow
\[
a:u\to p_P(x)
\]
such that
\[
R(a),
\qquad
O(\omega_P(a,x)),
\qquad
T(\omega_P(a,x),y),
\]
and
\[
\left(
\delta_P(a,x),
\sigma(\omega_P(a,x),y)
\right)
\models\varphi.
\]
\end{proposition}

\begin{proof}
\leavevmode
\begin{enumerate}
\item \textbf{Expand the restricted apex.}
The restricted one-step span has apex the subobject of \(M\) consisting of
triples \((a,x,y)\) satisfying the four pulled-back guards.

\item \textbf{Expand the diamond.}
The diamond says that there exists such a triple with source \((x,y)\) and
target satisfying \(\varphi\).

\item \textbf{Read the target.}
The target is
\[
\bigl(
\delta_P(a,x),
\sigma(\omega_P(a,x),y)
\bigr).
\]
\end{enumerate}
\end{proof}


\section{Guarded polynomial response liftings}
\label{sec:poly-relative-guarded-response-liftings}

\subsection{Existential profile liftings}
\label{subsec:poly-relative-existential-profile-liftings}

Assume \(\mathcal E\) is regular.  Let
\[
c_{\Phi,Y}:
S\to\mathcal M_{\mathbb A,\mathbb B,Y}(S)
\]
be the relative response coalgebra, and let
\[
\partial_S:\mathrm{Pos}_S\to\mathcal M_{\mathbb A,\mathbb B,Y}(S)
\]
be the position projection.

The position object carries evaluation maps
\[
\operatorname{in}:\mathrm{Pos}_S\to A_1,
\]
\[
\operatorname{out}:\mathrm{Pos}_S\to B_1,
\]
\[
\operatorname{down}:\mathrm{Pos}_S\to
\left(\int_{\mathbb B}Y\right)_1,
\]
and
\[
\operatorname{next}:\mathrm{Pos}_S\to S.
\]

The following liftings are coalgebraic predicate liftings for the response
polynomial, later pulled back along the coalgebra to the span modalities of
the flattened program category.  This is the same general pattern by which
modal operators for coalgebras are described using predicate liftings
\cite{Pattinson2003CoalgebraicModalLogic}.

\begin{definition}[Guarded existential profile lifting]
\label{def:poly-relative-guarded-existential-lifting}
Let
\[
R\hookrightarrow A_1,
\qquad
O\hookrightarrow B_1,
\qquad
T\hookrightarrow\left(\int_{\mathbb B}Y\right)_1,
\qquad
\varphi\hookrightarrow S
\]
be predicates.  The \textbf{guarded existential profile lifting} is the
predicate
\[
\lambda_{R/O/T}(\varphi)
\hookrightarrow
\mathcal M_{\mathbb A,\mathbb B,Y}(S)
\]
defined by
\[
\lambda_{R/O/T}(\varphi)
:=
\exists_{\partial_S}
\left(
\operatorname{in}^{\ast}R
\wedge
\operatorname{out}^{\ast}O
\wedge
\operatorname{down}^{\ast}T
\wedge
\operatorname{next}^{\ast}\varphi
\right).
\]
\end{definition}

Externally, a response profile satisfies
\[
\lambda_{R/O/T}(\varphi)
\]
when some input position satisfies \(R\), the selected raw output satisfies
\(O\), the downstream transition satisfies \(T\), and the continuation state
satisfies \(\varphi\).

\subsection{Agreement with span diamonds}
\label{subsec:poly-relative-lifting-span-agreement}

\begin{theorem}[Agreement of profile liftings and span diamonds]
\label{thm:poly-relative-lifting-span-agreement}
For every state guard
\[
\theta\hookrightarrow S,
\]
input guard
\[
R\hookrightarrow A_1,
\]
output guard
\[
O\hookrightarrow B_1,
\]
downstream transition guard
\[
T\hookrightarrow\left(\int_{\mathbb B}Y\right)_1,
\]
and postcondition
\[
\varphi\hookrightarrow S,
\]
one has
\[
\theta\wedge c_{\Phi,Y}^{\ast}\lambda_{R/O/T}(\varphi)
=
\langle\theta;R/O/T\rangle\varphi.
\]
\end{theorem}

\begin{proof}
\leavevmode
\begin{enumerate}
\item \textbf{Pull back the position projection.}
By Lemma~\ref{lem:poly-relative-mealy-position-pullback}, pulling back the
position projection
\[
\partial_S:\mathrm{Pos}_S\to\mathcal M_{\mathbb A,\mathbb B,Y}(S)
\]
along \(c_{\Phi,Y}\) gives the one-step execution object \(M\).

\item \textbf{Identify the evaluation maps.}
Under this identification,
\[
\operatorname{in}=\iota_A,
\qquad
\operatorname{out}=\iota_B,
\qquad
\operatorname{down}=\lambda_{\Phi,Y},
\qquad
\operatorname{next}=t_{\Phi,Y}.
\]

\item \textbf{Apply Beck--Chevalley.}
By Beck--Chevalley for regular images,
\[
c_{\Phi,Y}^{\ast}\lambda_{R/O/T}(\varphi)
=
\exists_{s_{\Phi,Y}}
\left(
\iota_A^\ast R
\wedge
\iota_B^\ast O
\wedge
\lambda_{\Phi,Y}^\ast T
\wedge
t_{\Phi,Y}^\ast\varphi
\right).
\]

\item \textbf{Pull in the state guard.}
Using Frobenius reciprocity,
\[
\theta\wedge c_{\Phi,Y}^{\ast}\lambda_{R/O/T}(\varphi)
=
\exists_{s_{\Phi,Y}}
\left(
s_{\Phi,Y}^\ast\theta
\wedge
\iota_A^\ast R
\wedge
\iota_B^\ast O
\wedge
\lambda_{\Phi,Y}^\ast T
\wedge
t_{\Phi,Y}^\ast\varphi
\right).
\]

\item \textbf{Identify the restricted diamond.}
The right-hand side is exactly the diamond of the restricted one-step span
\[
\theta;R/O/T.
\]
\end{enumerate}
\end{proof}

\begin{remark}
\label{rem:poly-relative-central-comparison}
The polynomial lifting is a predicate on total response profiles.  Pulling it
back along the relative Mealy coalgebra recovers the guarded span diamond on
the flattened program category.
\end{remark}

\section{Profile boxes}
\label{sec:poly-relative-profile-boxes}

\subsection{Weak and strong boxes}
\label{subsec:poly-relative-weak-strong-boxes}

Assume \(\mathcal E\) is a first-order Heyting category.

Let
\[
R\hookrightarrow A_1,
\qquad
O\hookrightarrow B_1,
\qquad
T\hookrightarrow\left(\int_{\mathbb B}Y\right)_1,
\qquad
\varphi\hookrightarrow S
\]
be predicates.

\begin{definition}[Weak universal profile box]
\label{def:poly-relative-weak-profile-box}
The \textbf{weak universal profile box} is
\[
\Box^{\mathrm{prof}}_{R/O/T}(\varphi)
:=
\forall_{\partial_S}
\left(
(\operatorname{in}^{\ast}R
\wedge
\operatorname{out}^{\ast}O
\wedge
\operatorname{down}^{\ast}T)
\Rightarrow
\operatorname{next}^{\ast}\varphi
\right).
\]
\end{definition}

Externally, this says that every position whose input, output, and downstream
transition satisfy the guards has a continuation satisfying \(\varphi\).

\begin{definition}[Strong universal profile box]
\label{def:poly-relative-strong-profile-box}
The \textbf{strong universal profile box} is
\[
\Box^{\mathrm{prof,strong}}_{R/O/T}(\varphi)
:=
\forall_{\partial_S}
\left(
\operatorname{in}^{\ast}R
\Rightarrow
\left(
\operatorname{out}^{\ast}O
\wedge
\operatorname{down}^{\ast}T
\wedge
\operatorname{next}^{\ast}\varphi
\right)
\right).
\]
\end{definition}

Externally, this says that every \(R\)-input position produces an output
satisfying \(O\), a downstream transition satisfying \(T\), and a continuation
satisfying \(\varphi\).

\subsection{Agreement with span boxes}
\label{subsec:poly-relative-profile-box-span-agreement}

\begin{theorem}[Agreement of weak profile boxes and span boxes]
\label{thm:poly-relative-profile-box-span-agreement}
One has
\[
c_{\Phi,Y}^{\ast}
\Box^{\mathrm{prof}}_{R/O/T}(\varphi)
=
[R/O/T]\varphi,
\]
where \([R/O/T]\varphi\) denotes the box modality of the restricted one-step
span determined by the guards \(R,O,T\).

More generally, for a state guard
\[
\theta\hookrightarrow S,
\]
one has
\[
[\theta;R/O/T]\varphi
=
\theta\Rightarrow
\bigl(
c_{\Phi,Y}^{\ast}
\Box^{\mathrm{prof}}_{R/O/T}(\varphi)
\bigr).
\]
\end{theorem}

\begin{proof}
\leavevmode
\begin{enumerate}
\item \textbf{Apply Beck--Chevalley for \(\forall\).}
By Beck--Chevalley for \(\forall\), pulling back
\(\Box^{\mathrm{prof}}_{R/O/T}(\varphi)\) along \(c_{\Phi,Y}\) gives
\[
\forall_{s_{\Phi,Y}}
\left(
(\iota_A^\ast R
\wedge
\iota_B^\ast O
\wedge
\lambda_{\Phi,Y}^\ast T)
\Rightarrow
t_{\Phi,Y}^\ast\varphi
\right).
\]

\item \textbf{Identify the restricted box.}
This is exactly the box of the restricted one-step span determined by
\(R,O,T\).

\item \textbf{Add the state guard.}
For the guarded version, use the test-prefix box law:
\[
[\theta?;U]\varphi
=
\theta\Rightarrow[U]\varphi.
\]
\end{enumerate}
\end{proof}

\begin{remark}[Weak versus strong profile boxes]
\label{rem:poly-relative-weak-strong-profile-boxes}
The weak profile box pulls back to the usual span box for the restricted
one-step span.  The strong profile box is generally stricter: it requires
every \(R\)-input position to produce an output satisfying \(O\), a downstream
transition satisfying \(T\), and a continuation satisfying \(\varphi\).
\end{remark}



\section{Sequential composition, finite choice, and conditionals}
\label{sec:poly-relative-tests-choice}

\subsection{Guarded commands}
\label{subsec:poly-relative-tests-guarded-commands}
Let
\[
U:S\nrightarrow S
\]
be an endoprogram.  The guarded command
\[
\theta?;U
\]
is the sequential composite of the test and \(U\).

\begin{proposition}[Guard-prefix law]
\label{prop:poly-relative-guard-prefix-law}
For every endoprogram \(U:S\nrightarrow S\),
\[
\langle\theta?;U\rangle\varphi
=
\theta\wedge\langle U\rangle\varphi.
\]
\end{proposition}

\begin{proof}
\leavevmode
\begin{enumerate}
\item \textbf{Use diamond functoriality.}
By functoriality of diamonds,
\[
\langle\theta?;U\rangle\varphi
=
\langle\theta?\rangle(\langle U\rangle\varphi).
\]

\item \textbf{Use the test law.}
The test law gives
\[
\langle\theta?\rangle\psi=\theta\wedge\psi.
\]

\item \textbf{Substitute.}
Substituting
\[
\psi=\langle U\rangle\varphi
\]
gives the claim.
\end{enumerate}
\end{proof}

\subsection{Finite guarded choice}
\label{subsec:poly-relative-finite-guarded-choice}

Assume now that \(\mathcal E\) is coherent.  Let
\[
T:S\nrightarrow S
\]
be a fixed ambient endospan, and let
\[
U,V\hookrightarrow T
\]
be subprograms.  If
\[
\theta,\epsilon\hookrightarrow S
\]
are state guards, then we identify
\[
\theta?;U
\]
with the subobject
\[
s_T^\ast\theta\wedge U\hookrightarrow T,
\]
and similarly identify
\[
\epsilon?;V
\]
with
\[
s_T^\ast\epsilon\wedge V\hookrightarrow T.
\]
Their guarded choice is the join of these subobjects in \(\mathrm{Sub}(T)\):
\[
(\theta?;U)\vee(\epsilon?;V)\hookrightarrow T.
\]

\begin{proposition}[Diamond law for guarded choice]
\label{prop:poly-relative-guarded-choice-diamond}
For guarded branches inside a common ambient endospan,
\[
\left\langle
(\theta?;U)\vee(\epsilon?;V)
\right\rangle\varphi
=
(\theta\wedge\langle U\rangle\varphi)
\vee
(\epsilon\wedge\langle V\rangle\varphi).
\]
\end{proposition}

\begin{proof}
\leavevmode
\begin{enumerate}
\item \textbf{Apply finite choice.}
By the finite choice law for subprograms in a coherent category,
\[
\left\langle
(\theta?;U)\vee(\epsilon?;V)
\right\rangle\varphi
=
\langle\theta?;U\rangle\varphi
\vee
\langle\epsilon?;V\rangle\varphi.
\]

\item \textbf{Apply the guard-prefix law.}
Using Proposition~\ref{prop:poly-relative-guard-prefix-law} on both branches
gives the displayed formula.
\end{enumerate}
\end{proof}

\subsection{Internal conditionals}
\label{subsec:poly-relative-conditionals}

\begin{definition}[Internal conditional]
\label{def:poly-relative-conditional}
Let
\[
T:S\nrightarrow S
\]
be a fixed ambient endospan, and let
\[
U,V\hookrightarrow T
\]
be subprograms.  Let
\[
\theta,\epsilon\hookrightarrow S
\]
be complementary guards, meaning
\[
\theta\wedge\epsilon=\bot,
\qquad
\theta\vee\epsilon=\top.
\]
The \textbf{internal conditional} is
\[
\mathbf{if}\ \theta\ \mathbf{then}\ U\ \mathbf{else}\ V
:=
(\theta?;U)\vee(\epsilon?;V)
\hookrightarrow T.
\]
In Boolean fibres one may take
\[
\epsilon=\neg\theta.
\]
\end{definition}


\section{Boxes for guarded choice}
\label{sec:poly-relative-boxes-choice}

\subsection{Boxes of tests and guarded commands}
\label{subsec:poly-relative-boxes-tests-guarded}

Assume \(\mathcal E\) is a first-order Heyting category.

For a test,
\[
[\theta?]\varphi=\theta\Rightarrow\varphi.
\]
Therefore
\[
[\theta?;U]\varphi
=
\theta\Rightarrow[U]\varphi.
\]

\subsection{Boxes of finite choice}
\label{subsec:poly-relative-boxes-finite-choice}

\begin{proposition}[Box of finite subprogram choice]
\label{prop:poly-relative-box-finite-choice}
For subprograms
\[
W,Z\hookrightarrow T
\]
inside a fixed ambient endospan \(T:S\nrightarrow S\),
\[
[W\vee Z]\varphi
=
[W]\varphi\wedge[Z]\varphi.
\]
\end{proposition}

\begin{proof}
\leavevmode
\begin{enumerate}
\item \textbf{Write the ambient span.}
Write
\[
T=(S\xleftarrow{\ell_T}T\xrightarrow{r_T}S).
\]

\item \textbf{Rewrite restricted boxes in the ambient fibre.}
For a restricted subspan \(i_W:W\hookrightarrow T\),
\[
[W]\varphi
=
\forall_{\ell_T i_W}\bigl((r_T i_W)^\ast\varphi\bigr).
\]
Equivalently,
\[
[W]\varphi
=
\forall_{\ell_T}\bigl(W\Rightarrow r_T^\ast\varphi\bigr),
\]
using the adjunction for the monomorphism \(i_W:W\hookrightarrow T\).

\item \textbf{Apply the same formula to the join.}
Hence
\[
[W\vee Z]\varphi
=
\forall_{\ell_T}
\bigl((W\vee Z)\Rightarrow r_T^\ast\varphi\bigr).
\]

\item \textbf{Use Heyting distributivity.}
In a Heyting algebra,
\[
(W\vee Z)\Rightarrow A
=
(W\Rightarrow A)\wedge(Z\Rightarrow A).
\]

\item \textbf{Use meet preservation by \(\forall\).}
Since \(\forall_{\ell_T}\) is a right adjoint, it preserves finite meets.
Therefore
\[
[W\vee Z]\varphi
=
[W]\varphi\wedge[Z]\varphi.
\]
\end{enumerate}
\end{proof}

\subsection{Boxes for conditionals}
\label{subsec:poly-relative-boxes-conditionals}

\begin{proposition}[Box law for conditionals]
\label{prop:poly-relative-conditional-box-law}
For complementary guards \(\theta,\epsilon\) and subprograms
\[
U,V\hookrightarrow T
\]
inside a common ambient endospan,
\[
\left[
\mathbf{if}\ \theta\ \mathbf{then}\ U\ \mathbf{else}\ V
\right]\varphi
=
(\theta\Rightarrow[U]\varphi)
\wedge
(\epsilon\Rightarrow[V]\varphi).
\]
\end{proposition}

\begin{proof}
\leavevmode
\begin{enumerate}
\item \textbf{Expand the conditional.}
The conditional is
\[
(\theta?;U)\vee(\epsilon?;V).
\]

\item \textbf{Apply the box law for finite choice.}
Proposition~\ref{prop:poly-relative-box-finite-choice} gives
\[
[(\theta?;U)\vee(\epsilon?;V)]\varphi
=
[\theta?;U]\varphi\wedge[\epsilon?;V]\varphi.
\]

\item \textbf{Use the test-prefix box law.}
The test-prefix box law gives
\[
[\theta?;U]\varphi=\theta\Rightarrow[U]\varphi
\]
and
\[
[\epsilon?;V]\varphi=\epsilon\Rightarrow[V]\varphi.
\]
\end{enumerate}
\end{proof}


\section{Domain and dynamic tests}
\label{sec:poly-relative-domain-tests}

\subsection{Domain predicates}
\label{subsec:poly-relative-domain-predicate}

Let
\[
U:S\nrightarrow S
\]
be an endoprogram.

\begin{definition}[Domain predicate]
\label{def:poly-relative-domain}
The \textbf{domain predicate} of \(U\) is
\[
\operatorname{dom}(U):=\langle U\rangle\top.
\]
\end{definition}

Thus \(\operatorname{dom}(U)\) is the predicate of states from which some
\(U\)-execution is possible.

\begin{proposition}[Domain-prefix law]
\label{prop:poly-relative-domain-prefix}
For every endoprogram \(U\),
\[
\langle \operatorname{dom}(U)?;U\rangle\varphi
=
\langle U\rangle\varphi.
\]
\end{proposition}

\begin{proof}
\leavevmode
\begin{enumerate}
\item \textbf{Apply the guard-prefix law.}
By the guard-prefix law,
\[
\langle \operatorname{dom}(U)?;U\rangle\varphi
=
\operatorname{dom}(U)\wedge\langle U\rangle\varphi.
\]

\item \textbf{Use monotonicity.}
Since
\[
\langle U\rangle\varphi\le\langle U\rangle\top
=
\operatorname{dom}(U),
\]
the meet is
\[
\langle U\rangle\varphi.
\]
\end{enumerate}
\end{proof}

\subsection{Antidomain predicates}
\label{subsec:poly-relative-antidomain-predicate}

If the predicate fibres are Boolean, define the antidomain predicate by
\[
\operatorname{adom}(U):=\neg\operatorname{dom}(U).
\]

\begin{proposition}[Antidomain annihilation]
\label{prop:poly-relative-antidomain-annihilation}
In Boolean predicate fibres,
\[
\langle \operatorname{adom}(U)?;U\rangle\varphi
=
\bot.
\]
\end{proposition}

\begin{proof}
\leavevmode
\begin{enumerate}
\item \textbf{Apply the guard-prefix law.}
By the guard-prefix law,
\[
\langle \operatorname{adom}(U)?;U\rangle\varphi
=
\operatorname{adom}(U)\wedge\langle U\rangle\varphi.
\]

\item \textbf{Use the domain bound.}
Since
\[
\langle U\rangle\varphi\le\operatorname{dom}(U)
\]
and
\[
\operatorname{adom}(U)=\neg\operatorname{dom}(U),
\]
the meet is bottom.
\end{enumerate}
\end{proof}

\section{Relative guarded Mealy automata}
\label{sec:poly-relative-guarded-automata}

\subsection{Guarded automata}
\label{subsec:poly-relative-guarded-automata-definition}

Let
\[
\Phi:\mathbb A\rightsquigarrow\mathbb B
\]
be a Mealy morphism and let \(Y\) be a right \(\mathbb B\)-action.

Let
\[
(\theta_i)_{i\in I}
\]
be a finite family of state guards on \(S_{\Phi,Y}\),
\[
(R_i)_{i\in I}
\]
a finite family of input-arrow guards on \(A_1\),
\[
(O_i)_{i\in I}
\]
a finite family of output-arrow guards on \(B_1\), and
\[
(T_i)_{i\in I}
\]
a finite family of downstream transition guards on
\[
\left(\int_{\mathbb B}Y\right)_1.
\]

Each quadruple determines a guarded one-step subprogram
\[
G_i:=(\theta_i;R_i/O_i/T_i)^\sharp\hookrightarrow M_{\Phi,Y}.
\]

\begin{definition}[Relative guarded Mealy automaton]
\label{def:poly-relative-guarded-mealy-automaton}
A \textbf{relative guarded Mealy automaton} over \((\Phi,Y)\) is a finite
family of guarded one-step subprograms
\[
G_i=(\theta_i;R_i/O_i/T_i)^\sharp\hookrightarrow M_{\Phi,Y}.
\]
Its one-step program is the finite join
\[
G:=\bigvee_{i\in I}G_i\hookrightarrow M_{\Phi,Y}.
\]
\end{definition}

\subsection{Modal semantics of guarded automata}
\label{subsec:poly-relative-guarded-automata-modal-semantics}

\begin{proposition}[Diamond semantics of guarded automata]
\label{prop:poly-relative-guarded-automaton-diamond}
For a relative guarded Mealy automaton with one-step program
\[
G=\bigvee_{i\in I}G_i,
\]
one has
\[
\langle G\rangle\varphi
=
\bigvee_{i\in I}\langle G_i\rangle\varphi.
\]
\end{proposition}

\begin{proof}
\leavevmode
\begin{enumerate}
\item \textbf{Use the common ambient object.}
All branches are subobjects of the common one-step execution object
\[
M_{\Phi,Y}.
\]

\item \textbf{Apply finite choice.}
The result is the finite choice law for subprograms inside this common ambient
object.
\end{enumerate}
\end{proof}

A relative guarded Mealy automaton is \textbf{witness-disjoint} if
\[
G_i\wedge G_j=\bot
\qquad
(i\ne j).
\]
It is \textbf{branch-disjoint} if
\[
\operatorname{dom}(G_i)\wedge\operatorname{dom}(G_j)=\bot
\qquad
(i\ne j).
\]
It is \textbf{guard-total} if
\[
\bigvee_{i\in I}\operatorname{dom}(G_i)=\top.
\]
When branch-disjointness and guard-totality hold, the branch domains form a
finite disjoint cover of the state object.  The enabled branch is unique at
the level of branch domains, though a branch may still contain multiple
execution witnesses.

\section{Output, downstream, and interface guards}
\label{sec:poly-relative-output-downstream-interface-guards}

\subsection{Output and downstream guards}
\label{subsec:poly-relative-output-downstream-guards}

The relative Mealy setting allows guards to inspect emitted output arrows and
the downstream transitions caused by those outputs.

Let
\[
U\hookrightarrow M_{\Phi,Y}
\]
be a one-step subprogram.

For an output guard
\[
O\hookrightarrow B_1,
\]
define
\[
U_O:=U\wedge\iota_B^\ast O\hookrightarrow M_{\Phi,Y}.
\]

For a downstream transition guard
\[
T\hookrightarrow\left(\int_{\mathbb B}Y\right)_1,
\]
define
\[
U_T:=U\wedge\lambda_{\Phi,Y}^\ast T\hookrightarrow M_{\Phi,Y}.
\]

\subsection{Interface guards}
\label{subsec:poly-relative-interface-guards}

Guards may also be placed directly on response interfaces.  Let
\[
\Gamma\hookrightarrow J_Y=A_0\times Y
\]
be an interface guard.  Pulling it back along
\[
j_{\Phi,Y}:S_{\Phi,Y}\to J_Y
\]
gives a state guard
\[
j_{\Phi,Y}^\ast\Gamma\hookrightarrow S_{\Phi,Y}.
\]
The pullback is displayed by
\[
\begin{tikzcd}[column sep=large,row sep=large]
j_{\Phi,Y}^\ast\Gamma
  \arrow[r]
  \arrow[d,hook]
&
\Gamma
  \arrow[d,hook]
\\
S_{\Phi,Y}
  \arrow[r,"j_{\Phi,Y}"']
&
J_Y
\arrow[ul,phantom,very near start,"\lrcorner"] .
\end{tikzcd}
\]

Thus one may form interface-guarded one-step programs such as
\[
(j_{\Phi,Y}^\ast\Gamma;R/O/T)^\sharp
\hookrightarrow M_{\Phi,Y}.
\]

\begin{proposition}[External reading of interface guards]
\label{prop:poly-relative-interface-guard-external}
In \(\mathbf{Set}\),
\[
(x,y)\models
\langle j_{\Phi,Y}^\ast\Gamma;R/O/T\rangle\varphi
\]
holds precisely when
\[
\Gamma(p_Px,y)
\]
and there exists an admissible input arrow
\[
a:u\to p_Px
\]
such that
\[
R(a),
\qquad
O(\omega_P(a,x)),
\qquad
T(\omega_P(a,x),y),
\]
and
\[
\left(
\delta_P(a,x),
\sigma(\omega_P(a,x),y)
\right)
\models\varphi.
\]
\end{proposition}

\begin{proof}
\leavevmode
\begin{enumerate}
\item \textbf{Identify the pulled-back state guard.}
The predicate
\[
j_{\Phi,Y}^\ast\Gamma
\]
holds at \((x,y)\) exactly when
\[
\Gamma(p_Px,y)
\]
holds.

\item \textbf{Apply the guarded one-step reading.}
The rest is Proposition~\ref{prop:poly-relative-guarded-primitive-external}
with
\[
\theta=j_{\Phi,Y}^\ast\Gamma.
\]
\end{enumerate}
\end{proof}


\section{The generated guarded program algebra}
\label{sec:poly-relative-generated-program-algebra}

\subsection{Generated programs}
\label{subsec:poly-relative-generated-programs}

\begin{definition}[Relative guarded Mealy program algebra]
\label{def:poly-relative-guarded-program-algebra}
The \textbf{relative guarded Mealy program algebra} generated by
\[
(\Phi,Y)
\]
is the class of endoprograms on
\[
S_{\Phi,Y}
\]
generated from:
\begin{enumerate}
\item tests \(\theta?\) for predicates
\[
\theta\hookrightarrow S_{\Phi,Y};
\]
\item guarded one-step programs
\[
\theta;R/O/T;
\]
\item sequential composition of spans;
\item finite guarded joins of subprograms inside a specified common ambient
span.
\end{enumerate}
\end{definition}

\subsection{Common ambient spans}
\label{subsec:poly-relative-common-ambient-spans}

\begin{remark}[Common ambient spans]
\label{rem:poly-relative-common-ambient-spans}
Finite guarded joins here are joins of subobjects inside a fixed ambient
execution span.  Choice between arbitrary generated spans requires coproducts
of span apexes and is a different, witnessful operation.
\end{remark}

\begin{remark}
\label{rem:poly-relative-program-algebra-from-polynomial}
The polynomial response profile supplies the curried one-step account of
guarded one-step programs.  Flattening returns the span calculus, where tests,
sequential composition, finite subprogram choice, diamonds, and boxes are
interpreted.
\end{remark}


\section{Stability under downstream change}
\label{subsec:relative-modal-stability-downstream}

Let
\[
f:Y\to Z
\]
be a morphism of right \(\mathbb B\)-actions. It induces
\[
h_f:S_{\Phi,Y}\to S_{\Phi,Z},
\qquad
h_f(x,y)=(x,f(y)),
\]
and
\[
k_f:M_{\Phi,Y}\to M_{\Phi,Z},
\qquad
k_f(a,x,y)=(a,x,f(y)).
\]
The source square is a pullback:
\[
\begin{tikzcd}
M_{\Phi,Y} \arrow[r,"k_f"] \arrow[d,"s_Y"']
&
M_{\Phi,Z} \arrow[d,"s_Z"]
\\
S_{\Phi,Y} \arrow[r,"h_f"']
&
S_{\Phi,Z}
\arrow[ul,phantom,very near start,"\lrcorner"]
\end{tikzcd}
.
\]
Indeed, both \(M_{\Phi,Y}\) and \(M_{\Phi,Z}\) are obtained by pulling back
\(t_A:A_1\to A_0\) along the respective structure maps to \(A_0\), and
\(h_f\) preserves those structure maps.

\begin{proposition}[Diamond stability under downstream change]
\label{prop:relative-diamond-stability-downstream}
Let \(\mathcal E\) be regular. For every subprogram
\[
V\hookrightarrow M_{\Phi,Z}
\]
and every predicate
\[
\psi\hookrightarrow S_{\Phi,Z},
\]
one has
\[
h_f^\ast\langle V\rangle\psi
=
\langle k_f^\ast V\rangle h_f^\ast\psi.
\]
\end{proposition}

\begin{proof}
Let
\[
i:V\hookrightarrow M_{\Phi,Z}
\]
be the subprogram inclusion, and form the pullback
\[
\begin{tikzcd}
k_f^\ast V \arrow[r,"\bar k_f"] \arrow[d,"\bar i"']
&
V \arrow[d,"i"]
\\
M_{\Phi,Y} \arrow[r,"k_f"']
&
M_{\Phi,Z}
\arrow[ul,phantom,very near start,"\lrcorner"]
\end{tikzcd}
.
\]
Because the source square
\[
\begin{tikzcd}
M_{\Phi,Y} \arrow[r,"k_f"] \arrow[d,"s_Y"']
&
M_{\Phi,Z} \arrow[d,"s_Z"]
\\
S_{\Phi,Y} \arrow[r,"h_f"']
&
S_{\Phi,Z}
\end{tikzcd}
\]
is a pullback, the induced square
\[
\begin{tikzcd}
k_f^\ast V \arrow[r,"\bar k_f"] \arrow[d,"s_Y\bar i"']
&
V \arrow[d,"s_Z i"]
\\
S_{\Phi,Y} \arrow[r,"h_f"']
&
S_{\Phi,Z}
\end{tikzcd}
\]
is also a pullback. Moreover target compatibility gives
\[
t_Z i\bar k_f=h_f t_Y\bar i.
\]
Therefore Beck--Chevalley for regular images gives
\[
h_f^\ast\exists_{s_Z i}\bigl((t_Z i)^\ast\psi\bigr)
=
\exists_{s_Y\bar i}\bar k_f^\ast\bigl((t_Z i)^\ast\psi\bigr).
\]
Using target compatibility, the right-hand side becomes
\[
\exists_{s_Y\bar i}\bigl((t_Y\bar i)^\ast h_f^\ast\psi\bigr)
=
\langle k_f^\ast V\rangle h_f^\ast\psi.
\]
\end{proof}

\begin{proposition}[Box stability under downstream change]
\label{prop:relative-box-stability-downstream}
Let \(\mathcal E\) be first-order Heyting. For every subprogram
\[
V\hookrightarrow M_{\Phi,Z}
\]
and every predicate
\[
\psi\hookrightarrow S_{\Phi,Z},
\]
one has
\[
h_f^\ast[V]\psi
=
[k_f^\ast V]h_f^\ast\psi.
\]
\end{proposition}

\begin{proof}
With notation as in the proof of
Proposition~\ref{prop:relative-diamond-stability-downstream}, the square
\[
\begin{tikzcd}
k_f^\ast V \arrow[r,"\bar k_f"] \arrow[d,"s_Y\bar i"']
&
V \arrow[d,"s_Z i"]
\\
S_{\Phi,Y} \arrow[r,"h_f"']
&
S_{\Phi,Z}
\end{tikzcd}
\]
is a pullback, and target compatibility gives
\[
t_Z i\bar k_f=h_f t_Y\bar i.
\]
Hence Beck--Chevalley for the right adjoints \(\forall\) gives
\[
h_f^\ast\forall_{s_Z i}\bigl((t_Z i)^\ast\psi\bigr)
=
\forall_{s_Y\bar i}\bar k_f^\ast\bigl((t_Z i)^\ast\psi\bigr).
\]
Using target compatibility, this is
\[
\forall_{s_Y\bar i}\bigl((t_Y\bar i)^\ast h_f^\ast\psi\bigr)
=
[k_f^\ast V]h_f^\ast\psi.
\]
\end{proof}


\section{Simulation stability after flattening}
\label{sec:poly-relative-simulation-stability}

\subsection{Guard transport along simulations}
\label{subsec:poly-relative-guard-transport-simulation}

Let
\[
\Theta:\Phi\Rightarrow\Psi
\]
be a Mealy simulation, and let \(Y\) be a right \(\mathbb B\)-action.  The
induced functor
\[
H_{\Theta,Y}:
\mathsf{Prog}_{\Psi,Y}\to\mathsf{Prog}_{\Phi,Y}
\]
has object map
\[
h_{\Theta,Y}:S_{\Psi,Y}\to S_{\Phi,Y}
\]
and arrow map
\[
k_{\Theta,Y}:M_{\Psi,Y}\to M_{\Phi,Y},
\]
given by
\[
k_{\Theta,Y}(a,y,z)
=
\bigl(a,\theta y,\sigma(\alpha_y,z)\bigr).
\]

The source square is a pullback:
\[
\begin{tikzcd}[column sep=large,row sep=large]
M_{\Psi,Y}
  \arrow[r,"k_{\Theta,Y}"]
  \arrow[d,"s_\Psi"']
&
M_{\Phi,Y}
  \arrow[d,"s_\Phi"]
\\
S_{\Psi,Y}
  \arrow[r,"h_{\Theta,Y}"']
&
S_{\Phi,Y}
\arrow[ul,phantom,very near start,"\lrcorner"] .
\end{tikzcd}
\]

\begin{proposition}[Guard transport along simulations]
\label{prop:poly-relative-guard-transport-simulation}
Let
\[
U\hookrightarrow M_{\Phi,Y}
\]
be a one-step subprogram for \(\Phi\).  Its pullback along the simulation is
\[
k_{\Theta,Y}^\ast U\hookrightarrow M_{\Psi,Y}.
\]
If \(U\) is presented by guards \(R,O,T\) on the target side, then
\(k_{\Theta,Y}^\ast U\) is the source-side predicate requiring
\[
R(a),
\]
\[
O(\omega_P(a,\theta y)),
\]
and
\[
T(\omega_P(a,\theta y),\sigma(\alpha_y,z)).
\]
If the target-side subprogram also has a state guard
\[
\theta_0\hookrightarrow S_{\Phi,Y},
\]
then its pullback additionally requires
\[
\theta_0(\theta y,\sigma(\alpha_y,z)).
\]
Equivalently, the state guard is pulled back as
\[
h_{\Theta,Y}^{\ast}\theta_0.
\]
\end{proposition}

\begin{proof}
\leavevmode
\begin{enumerate}
\item \textbf{Track the input component.}
The map \(k_{\Theta,Y}\) preserves the input component \(a\).  Hence input
guards pull back by direct substitution:
\[
R(a).
\]

\item \textbf{Track the target output component.}
The output and downstream transition components are those of the
\(\Phi\)-profile at the evaluated state
\[
(\theta y,\sigma(\alpha_y,z)).
\]
Thus the output guard becomes
\[
O(\omega_P(a,\theta y)).
\]

\item \textbf{Track the downstream transition.}
The downstream transition guard becomes
\[
T(\omega_P(a,\theta y),\sigma(\alpha_y,z)).
\]

\item \textbf{Track state guards.}
A state guard is pulled back along the object map \(h_{\Theta,Y}\), hence
becomes
\[
h_{\Theta,Y}^{\ast}\theta_0.
\]
\end{enumerate}
\end{proof}

\subsection{Modal stability under simulations}
\label{subsec:poly-relative-modal-stability-simulation}

\begin{proposition}[Diamond stability under simulations]
\label{prop:poly-relative-diamond-stability-simulation}
Assume \(\mathcal E\) is regular.  For every one-step subprogram
\[
U\hookrightarrow M_{\Phi,Y}
\]
and every predicate
\[
\varphi\hookrightarrow S_{\Phi,Y},
\]
one has
\[
h_{\Theta,Y}^\ast\langle U\rangle\varphi
=
\langle k_{\Theta,Y}^\ast U\rangle h_{\Theta,Y}^\ast\varphi.
\]
\end{proposition}

\begin{proof}
\leavevmode
\begin{enumerate}
\item \textbf{Pull back the subprogram.}
Let
\[
i:U\hookrightarrow M_{\Phi,Y}
\]
be the subprogram inclusion, and form the pullback
\[
\begin{tikzcd}[column sep=large,row sep=large]
k_{\Theta,Y}^\ast U
  \arrow[r,"\bar k_{\Theta,Y}"]
  \arrow[d,"\bar i"']
&
U
  \arrow[d,"i"]
\\
M_{\Psi,Y}
  \arrow[r,"k_{\Theta,Y}"']
&
M_{\Phi,Y}
\arrow[ul,phantom,very near start,"\lrcorner"] .
\end{tikzcd}
\]

\item \textbf{Use the source pullback square.}
Since the square
\[
\begin{tikzcd}[column sep=large,row sep=large]
M_{\Psi,Y} \arrow[r,"k_{\Theta,Y}"] \arrow[d,"s_\Psi"']
&
M_{\Phi,Y} \arrow[d,"s_\Phi"]
\\
S_{\Psi,Y} \arrow[r,"h_{\Theta,Y}"']
&
S_{\Phi,Y}
\end{tikzcd}
\]
is a pullback, the induced square
\[
\begin{tikzcd}[column sep=large,row sep=large]
k_{\Theta,Y}^\ast U
  \arrow[r,"\bar k_{\Theta,Y}"]
  \arrow[d,"s_\Psi\bar i"']
&
U
  \arrow[d,"s_\Phi i"]
\\
S_{\Psi,Y}
  \arrow[r,"h_{\Theta,Y}"']
&
S_{\Phi,Y}
\arrow[ul,phantom,very near start,"\lrcorner"]
\end{tikzcd}
\]
is also a pullback.

\item \textbf{Use target compatibility.}
Since \(H_{\Theta,Y}\) is an internal functor,
\[
t_\Phi i\bar k_{\Theta,Y}=h_{\Theta,Y}t_\Psi\bar i.
\]

\item \textbf{Apply Beck--Chevalley.}
Beck--Chevalley for regular images gives
\[
h_{\Theta,Y}^\ast
\exists_{s_\Phi i}\bigl((t_\Phi i)^\ast\varphi\bigr)
=
\exists_{s_\Psi\bar i}
\bigl((t_\Psi\bar i)^\ast h_{\Theta,Y}^\ast\varphi\bigr).
\]

\item \textbf{Identify the two diamonds.}
The left-hand side is
\[
h_{\Theta,Y}^\ast\langle U\rangle\varphi,
\]
and the right-hand side is
\[
\langle k_{\Theta,Y}^\ast U\rangle h_{\Theta,Y}^\ast\varphi.
\]
\end{enumerate}
\end{proof}

\begin{proposition}[Box stability under simulations]
\label{prop:poly-relative-box-stability-simulation}
Assume \(\mathcal E\) is first-order Heyting.  For every one-step subprogram
\[
U\hookrightarrow M_{\Phi,Y}
\]
and every predicate
\[
\varphi\hookrightarrow S_{\Phi,Y},
\]
one has
\[
h_{\Theta,Y}^\ast[U]\varphi
=
[k_{\Theta,Y}^\ast U]h_{\Theta,Y}^\ast\varphi.
\]
\end{proposition}

\begin{proof}
\leavevmode
\begin{enumerate}
\item \textbf{Use the same pullback square.}
Use the pullback square for
\[
k_{\Theta,Y}^\ast U
\]
constructed in the proof of
Proposition~\ref{prop:poly-relative-diamond-stability-simulation}.

\item \textbf{Use target compatibility.}
Again,
\[
t_\Phi i\bar k_{\Theta,Y}=h_{\Theta,Y}t_\Psi\bar i.
\]

\item \textbf{Apply Beck--Chevalley for \(\forall\).}
Beck--Chevalley for the right adjoints \(\forall\) gives the displayed
equality.
\end{enumerate}
\end{proof}

\begin{remark}
\label{rem:poly-relative-simulation-stability-meaning}
Because simulations are output-lax, guard transport is not obtained by
substituting the source output \(\omega_Q\) for the target output
\(\omega_P\).  It is obtained by pulling the target subprogram back along the
actual flattened functor \(H_{\Theta,Y}\).
\end{remark}



\section{Finite-trajectory modalities}
\label{sec:dynamic-finite-trajectory-modalities}

The one-step program category is an internal category, so its finite paths
supply a family of execution spans.  This section applies the span
predicate-transformer calculus to those path objects.

\subsection{Trajectory spans and endpoint modalities}
\label{subsec:dynamic-trajectory-endpoint-modalities}

For \(n\geq0\), write
\[
    \mathscr T_{\Phi,Y}(n)
    :=
    \operatorname{Path}(\mathsf{Prog}_{\Phi,Y})(n).
\]
Let
\[
    \nu_i^{(n)}:
    \mathscr T_{\Phi,Y}(n)\to S_{\Phi,Y},
    \qquad
    0\leq i\leq n,
\]
be the \(i\)-th boundary-state map.  In particular,
\[
    s_n:=\nu_0^{(n)},
    \qquad
    t_n:=\nu_n^{(n)}.
\]
The length-\(n\) trajectory span is
\[
    \mathbf t_{\Phi,Y}^{(n)}
    :=
    \left(
    S_{\Phi,Y}
    \xleftarrow{s_n}
    \mathscr T_{\Phi,Y}(n)
    \xrightarrow{t_n}
    S_{\Phi,Y}
    \right).
\]

\begin{definition}[Finite-horizon endpoint modalities]
\label{def:dynamic-finite-horizon-endpoint-modalities}
Assume \(\mathcal E\) is regular.  For a predicate
\[
    \varphi\hookrightarrow S_{\Phi,Y},
\]
define
\[
    \langle n\rangle_{\Phi,Y}\varphi
    :=
    \exists_{s_n}(t_n^\ast\varphi).
\]
If \(\mathcal E\) is first-order Heyting, define
\[
    [n]_{\Phi,Y}\varphi
    :=
    \forall_{s_n}(t_n^\ast\varphi).
\]
\end{definition}

Externally,
\[
    x\models\langle n\rangle_{\Phi,Y}\varphi
\]
means that some \(n\)-segment relative execution begins at \(x\) and ends in
\(\varphi\).  The formula
\[
    x\models[n]_{\Phi,Y}\varphi
\]
means that every such \(n\)-segment execution ends in \(\varphi\).

At \(n=0\), the only trajectory is the identity trajectory, so
\[
    \langle0\rangle_{\Phi,Y}\varphi=\varphi,
    \qquad
    [0]_{\Phi,Y}\varphi=\varphi.
\]

\subsection{Concatenation laws}
\label{subsec:dynamic-trajectory-concatenation-laws}

The Segal isomorphism identifies a path of length \(m+n\) with a path of
length \(m\) followed by a path of length \(n\):
\[
    \mathscr T_{\Phi,Y}(m+n)
    \cong
    \mathscr T_{\Phi,Y}(m)
    \times_{S_{\Phi,Y}}
    \mathscr T_{\Phi,Y}(n),
\]
where the pullback matches the final state of the first path with the initial
state of the second.

\begin{proposition}[Composition of finite-horizon modalities]
\label{prop:dynamic-finite-horizon-composition}
For \(m,n\geq0\), one has
\[
    \langle m+n\rangle_{\Phi,Y}
    =
    \langle m\rangle_{\Phi,Y}
    \circ
    \langle n\rangle_{\Phi,Y}.
\]
If \(\mathcal E\) is first-order Heyting, then
\[
    [m+n]_{\Phi,Y}
    =
    [m]_{\Phi,Y}
    \circ
    [n]_{\Phi,Y}.
\]
\end{proposition}

\begin{proof}
Under the displayed Segal isomorphism, the length-\((m+n)\) trajectory span
is the composite of the length-\(m\) and length-\(n\) trajectory spans.
The diamond identity is therefore
Proposition~\ref{prop:relative-diamond-functoriality}, and the box identity is
Proposition~\ref{prop:relative-box-functoriality}.
\end{proof}

\subsection{Boundary modalities}
\label{subsec:dynamic-trajectory-boundary-modalities}

Endpoint modalities see only the last boundary.  To express invariants and
safety conditions, one must also inspect intermediate boundaries.

\begin{definition}[Boundary modalities]
\label{def:dynamic-trajectory-boundary-modalities}
For \(0\leq i\leq n\), define
\[
    \langle n,i\rangle_{\Phi,Y}\varphi
    :=
    \exists_{s_n}
    \left(
        \bigl(\nu_i^{(n)}\bigr)^\ast\varphi
    \right).
\]
When \(\forall_{s_n}\) exists, define
\[
    [n,i]_{\Phi,Y}\varphi
    :=
    \forall_{s_n}
    \left(
        \bigl(\nu_i^{(n)}\bigr)^\ast\varphi
    \right).
\]
\end{definition}

Thus \(\langle n,i\rangle\varphi\) asserts that some length-\(n\) execution
reaches \(\varphi\) at boundary \(i\), while \([n,i]\varphi\) asserts that all
length-\(n\) executions lie in \(\varphi\) at that boundary.

For \(i=n\),
\[
    \langle n,n\rangle_{\Phi,Y}
    =
    \langle n\rangle_{\Phi,Y},
    \qquad
    [n,n]_{\Phi,Y}
    =
    [n]_{\Phi,Y}.
\]
For \(i=0\), the boundary is the initial state.  Since every state has an
all-identity path of every finite length, the source map \(s_n\) has a section,
and hence
\[
    \langle n,0\rangle_{\Phi,Y}\varphi=\varphi,
    \qquad
    [n,0]_{\Phi,Y}\varphi=\varphi.
\]

\begin{definition}[Finite-horizon safety transformer]
\label{def:dynamic-finite-horizon-safety}
For a predicate \(C\hookrightarrow S_{\Phi,Y}\), define
\[
    \operatorname{Safe}_n(C)
    :=
    \bigwedge_{0\leq i\leq n}
    [n,i]_{\Phi,Y}C.
\]
When the required small meets exist, define
\[
    \operatorname{Safe}_{\mathrm{fin}}(C)
    :=
    \bigwedge_{n\in\mathbb N}
    \operatorname{Safe}_n(C).
\]
\end{definition}

\subsection{Guards on complete trajectories}
\label{subsec:dynamic-guarded-trajectory-modalities}

A trajectory may be restricted before applying a modality.  Let
\[
    G\hookrightarrow\mathscr T_{\Phi,Y}(n)
\]
be any trajectory guard, with inclusion \(i_G\).  Define
\[
    \langle G\rangle\varphi
    :=
    \exists_{s_ni_G}
    \left(
        (t_ni_G)^\ast\varphi
    \right),
\]
and, when available,
\[
    [G]\varphi
    :=
    \forall_{s_ni_G}
    \left(
        (t_ni_G)^\ast\varphi
    \right).
\]
The guard \(G\) may be obtained by pulling back predicates along:

\begin{enumerate}
\item the complete input-path label;
\item the complete output-path label;
\item the downstream trajectory label;
\item any boundary-state map \(\nu_i^{(n)}\);
\item any segment projection to the one-step execution object.
\end{enumerate}

Consequently, trajectory formulas can constrain both the endpoint and the
internal segmentation of an execution.

\subsection{Contraction and the limits of horizon counting}
\label{subsec:dynamic-trajectory-contraction-collapse}

For every \(n\geq1\), categorical composition gives a contraction map
\[
    \mu_n:
    \mathscr T_{\Phi,Y}(n)\to M_{\Phi,Y}
\]
satisfying
\[
    s_{\Phi,Y}\mu_n=s_n,
    \qquad
    t_{\Phi,Y}\mu_n=t_n.
\]
Conversely, insertion of identities gives a segmentation map
\[
    \eta_n:
    M_{\Phi,Y}\to\mathscr T_{\Phi,Y}(n)
\]
with the same source and target and
\[
    \mu_n\eta_n=1_{M_{\Phi,Y}}.
\]

\begin{proposition}[Collapse of unrestricted endpoint horizons]
\label{prop:dynamic-unrestricted-endpoint-collapse}
For every \(n\geq1\) and every predicate
\[
    \varphi\hookrightarrow S_{\Phi,Y},
\]
one has
\[
    \langle n\rangle_{\Phi,Y}\varphi
    =
    \langle\mathbf m_{\Phi,Y}\rangle\varphi.
\]
If \(\mathcal E\) is first-order Heyting, then
\[
    [n]_{\Phi,Y}\varphi
    =
    [\mathbf m_{\Phi,Y}]\varphi.
\]
\end{proposition}

\begin{proof}
The contraction map sends every length-\(n\) trajectory to a one-step
categorical execution with the same source and target.  This gives one
inclusion for both modalities.  The identity-segmentation map sends every
one-step execution to a length-\(n\) trajectory with the same source and
target, giving the reverse inclusion.  Equivalently, the two spans have the
same induced endpoint relation, and the maps \(\mu_n\) and \(\eta_n\) witness
both comparisons.
\end{proof}

\begin{remark}[Segmentation-sensitive formulas]
\label{rem:dynamic-segmentation-sensitive-formulas}
The preceding collapse applies only to unrestricted endpoint properties.
It does not identify:

\begin{enumerate}
\item boundary modalities \([n,i]\) for \(0<i<n\);
\item trajectory guards that inspect individual segments;
\item paths restricted to a chosen generating graph;
\item letterwise paths inside a free word category.
\end{enumerate}

Thus categorical path length measures a chosen segmentation.  A genuinely
atomic temporal semantics requires a specified class of generators.
\end{remark}


\section{Residual reachability and universal finite-path safety}
\label{sec:dynamic-residual-reachability-safety}

Chapter~\ref{ch:temporal-trajectories-segal-execution} proved that generated
derivative closure is endpoint reachability along finite canonical residual
executions.  We now express the right adjoint of derivative invariance as a
trajectory box.  Write
\[
    \mathbf{Beh}
    =
    \mathbf{Beh}^{\mathrm{can}}_{\mathbb A,\mathbb B}
\]
for the canonical behaviour machine, and let
\[
    \mathscr T_{\mathbf{Beh}}(n)
\]
be its length-\(n\) trajectory object.  The notation \(s_n,t_n\) and
\(\nu_i^{(n)}\) is as in
Section~\ref{sec:dynamic-finite-trajectory-modalities}.
\subsection{The derivative safety core as universal path safety}
\label{subsec:temporal-interpretation-safety-core}

For a behavioural predicate
\[
    C\hookrightarrow\mathbf{Beh}_0,
\]
define
\[
    \operatorname{Safe}_{\mathrm{path}}(C)
    :=
    \bigwedge_{n\in\mathbb N}
    \;
    \bigwedge_{0\leq i\leq n}
    \forall_{s_n}
    \left(
        \bigl(\nu_i^{(n)}\bigr)^\ast C
    \right)
    \hookrightarrow
    \mathbf{Beh}_0.
\]
Here
\[
    \forall_{s_n}:
    \operatorname{Sub}
    \bigl(
        \mathscr T_{\mathbf{Beh}}(n)
    \bigr)
    \to
    \operatorname{Sub}(\mathbf{Beh}_0)
\]
is the dependent-product right adjoint to pullback along \(s_n\).

By adjunction, a subobject
\[
    S\hookrightarrow\mathbf{Beh}_0
\]
satisfies
\[
    S\leq\operatorname{Safe}_{\mathrm{path}}(C)
\]
if and only if
\[
    s_n^\ast S
    \leq
    \bigl(\nu_i^{(n)}\bigr)^\ast C
\]
for every \(n\) and every boundary \(i\). Equivalently, every finite
canonical execution beginning in \(S\) remains in \(C\) at every boundary.
Since each intermediate boundary is the endpoint of a prefix, one may also
write
\[
    \operatorname{Safe}_{\mathrm{path}}(C)
    =
    \bigwedge_{n\in\mathbb N}
    \forall_{s_n}(t_n^\ast C).
\]

For \(n\geq0\) and \(0\leq i\leq n\), let
\[
    \kappa_{n,i}:
    \mathscr T_{\mathbf{Beh}}(n)
    \to
    \mathscr T_{\mathbf{Beh}}(1)
\]
be contraction of the prefix from boundary \(0\) to boundary \(i\). It is
induced by
\[
    [1]\to[n],
    \qquad
    0\mapsto0,
    \quad
    1\mapsto i.
\]
When \(i=0\), the map inserts the identity at the initial boundary. It
satisfies
\[
    s_1\kappa_{n,i}=s_n,
    \qquad
    t_1\kappa_{n,i}=\nu_i^{(n)}.
\]

\begin{theorem}[Safety core as universal finite-path safety]
\label{thm:safety-core-universal-temporal-safety}
For every behavioural predicate \(C\),
\[
    \operatorname{Safe}_{\mathrm{path}}(C)
    =
    \operatorname{Safe}_{\partial}(C).
\]
Consequently, \(\operatorname{Safe}_{\partial}(C)\) is the largest subobject
\(S\hookrightarrow\mathbf{Beh}_0\) from which every finite canonical
residual execution remains in \(C\) at every boundary.
\end{theorem}

\begin{proof}
Put
\[
    Q_C
    :=
    \forall_{s_1}(t_1^\ast C).
\]
By definition,
\[
    s_1^\ast Q_C
    \leq
    t_1^\ast C.
\]
Pulling this inequality back along \(\kappa_{n,i}\) gives
\[
    s_n^\ast Q_C
    \leq
    \bigl(\nu_i^{(n)}\bigr)^\ast C.
\]
By adjunction,
\[
    Q_C
    \leq
    \forall_{s_n}
    \left(
        \bigl(\nu_i^{(n)}\bigr)^\ast C
    \right)
\]
for every \(n\) and \(i\). Hence
\[
    Q_C
    \leq
    \operatorname{Safe}_{\mathrm{path}}(C).
\]
Conversely, the defining meet for
\(\operatorname{Safe}_{\mathrm{path}}(C)\) contains the factor corresponding
to \(n=1\) and \(i=1\), so
\[
    \operatorname{Safe}_{\mathrm{path}}(C)
    \leq
    Q_C.
\]
Therefore
\[
    \operatorname{Safe}_{\mathrm{path}}(C)=Q_C.
\]

The predicate \(Q_C\) is contained in \(C\), because identity executions are
among the one-step executions. It is derivative invariant. Indeed, suppose
\(H\in Q_C\), let \(a\) be applicable to \(H\), and let \(b\) be applicable
to \(\partial_aH\). The state
\[
    \partial_b(\partial_aH)
\]
is the endpoint of the one-step execution from \(H\) labelled by the
composite input arrow. Since \(H\in Q_C\), this endpoint lies in \(C\).
Thus every derivative of \(\partial_aH\) lies in \(C\), so
\(\partial_aH\in Q_C\).

Hence \(Q_C\) is derivative invariant and lies below \(C\). Conversely, any
derivative-invariant subobject below \(C\) satisfies the one-step condition
and therefore lies below \(Q_C\). By the universal property of the derivative
safety core,
\[
    Q_C=\operatorname{Safe}_{\partial}(C).
\]
Combining the two equalities proves the theorem.
\end{proof}

\begin{corollary}[Finite-path meaning of the derivative adjoint triple]
\label{cor:temporal-meaning-derivative-adjoint-triple}
The adjoint triple
\[
    \operatorname{Gen}_{\partial}
    \dashv
    I_{\partial}
    \dashv
    \operatorname{Safe}_{\partial}
\]
has the following execution-theoretic interpretation:
\begin{enumerate}
    \item \(\operatorname{Gen}_{\partial}(C)\) consists of residual states
    reachable as endpoints of finite canonical executions beginning in \(C\);
    \item \(I_{\partial}\) forgets derivative invariance and retains the
    underlying behavioural predicate;
    \item \(\operatorname{Safe}_{\partial}(C)\) consists of residual states
    from which every finite canonical execution remains in \(C\).
\end{enumerate}
\end{corollary}



\section{Modal invariance under quotients and minimization}
\label{sec:dynamic-modal-invariance-quotients}

The quotient theorems of
Chapter~\ref{ch:temporal-trajectories-segal-execution} are structural: they
produce regular epimorphisms on states, one-step executions, and every finite
trajectory object.  Logical invariance requires one further condition: the
predicates and guards under discussion must descend to the quotient.

\subsection{Cartesian comparison of execution spans}
\label{subsec:dynamic-cartesian-span-comparison}

Consider spans
\[
    R=
    \left(
    X\xleftarrow{\ell_R}R\xrightarrow{r_R}X
    \right),
    \qquad
    \overline R=
    \left(
    Q\xleftarrow{\ell_{\overline R}}\overline R
    \xrightarrow{r_{\overline R}}Q
    \right),
\]
and maps
\[
    e:X\to Q,
    \qquad
    k:R\to\overline R
\]
such that
\[
    e\ell_R=\ell_{\overline R}k,
    \qquad
    er_R=r_{\overline R}k.
\]

\begin{definition}[Source-cartesian span comparison]
\label{def:dynamic-source-cartesian-span-comparison}
The pair \((e,k)\) is \textbf{source-cartesian} if
\[
\begin{tikzcd}
R \arrow[r,"k"] \arrow[d,"\ell_R"']
&
\overline R \arrow[d,"\ell_{\overline R}"]
\\
X \arrow[r,"e"']
&
Q
\arrow[ul,phantom,very near start,"\lrcorner"]
\end{tikzcd}
\]
is a pullback.
\end{definition}

\begin{proposition}[Modal Beck--Chevalley for span comparisons]
\label{prop:dynamic-modal-beck-chevalley-span-comparison}
If \((e,k)\) is source-cartesian and \(\mathcal E\) is regular, then
\[
    e^\ast\langle\overline R\rangle\varphi
    =
    \langle R\rangle e^\ast\varphi.
\]
If \(\mathcal E\) is first-order Heyting, then
\[
    e^\ast[\overline R]\varphi
    =
    [R]e^\ast\varphi.
\]
\end{proposition}

\begin{proof}
For diamonds, apply Beck--Chevalley for regular images to the source
pullback square and use target compatibility:
\[
\begin{aligned}
e^\ast\exists_{\ell_{\overline R}}
    (r_{\overline R}^\ast\varphi)
&=
\exists_{\ell_R}
    k^\ast(r_{\overline R}^\ast\varphi)\\
&=
\exists_{\ell_R}
    (r_R^\ast e^\ast\varphi).
\end{aligned}
\]
The proof for boxes is identical with \(\forall\) in place of \(\exists\).
\end{proof}

\subsection{Strict regular quotients}
\label{subsec:dynamic-modal-strict-regular-quotients}

Let
\[
    e:\Phi\twoheadrightarrow\Psi
\]
be a strict regular quotient over fixed input and output interfaces.  In every
finite degree, deterministic path lifting identifies
\[
    \mathscr T_\Phi(n)
    \cong
    P\times_{A_0}\operatorname{Path}(\mathbb A^{\mathrm{op}})(n)
\]
and similarly for \(\Psi\).  Therefore the square formed by the initial-state
maps and
\[
    \mathscr T_e(n):
    \mathscr T_\Phi(n)\to\mathscr T_\Psi(n)
\]
is a pullback.

\begin{theorem}[Modal invariance of descended trajectory predicates]
\label{thm:dynamic-modal-invariance-trajectory-quotients}
Let
\[
    \overline\varphi\hookrightarrow Q
\]
be a predicate on the quotient state object and put
\[
    \varphi:=e^\ast\overline\varphi.
\]
Then, for every \(n\geq0\),
\[
    e^\ast
    \langle n\rangle_{\Psi}\overline\varphi
    =
    \langle n\rangle_{\Phi}\varphi.
\]
If \(\mathcal E\) is first-order Heyting, then
\[
    e^\ast
    [n]_{\Psi}\overline\varphi
    =
    [n]_{\Phi}\varphi.
\]
The same statements hold for boundary modalities after replacing the target
map by the corresponding boundary-state map.
\end{theorem}

\begin{proof}
The source square is a pullback by deterministic path lifting, and every
boundary-state map commutes with \(\mathscr T_e(n)\).  Apply
Proposition~\ref{prop:dynamic-modal-beck-chevalley-span-comparison}.
\end{proof}

\subsection{Guards and behavioural minimization}
\label{subsec:dynamic-guards-behavioural-minimization}

A state predicate, input guard, output guard, downstream guard, or trajectory
guard is \textbf{quotient-saturated} when it is the inverse image of a
predicate on the corresponding quotient object.  For quotient-saturated
guards, the restricted execution squares remain pullbacks, so the same
Beck--Chevalley argument applies.

\begin{corollary}[Behavioural invariance of descended formulas]
\label{cor:dynamic-behavioural-invariance-descended-formulas}
Let
\[
    e:P\twoheadrightarrow\operatorname{Beh}(P)
\]
be the behavioural quotient, under the hypotheses of the behavioural
minimization theorem.  Every formula built from quotient-saturated atomic
predicates and guards using:

\begin{enumerate}
\item finite meets, finite joins, and Heyting implication when available;
\item tests;
\item guarded one-step diamonds and boxes;
\item finite guarded choice;
\item finite-trajectory endpoint and boundary modalities;
\end{enumerate}

is invariant under \(e\): its interpretation on \(P\) is the inverse image of
its interpretation on \(\operatorname{Beh}(P)\).
\end{corollary}

\begin{proof}
Proceed by induction on formulas and generated programs.  Pullback preserves
finite meets and, in the coherent setting, the finite joins in question.
Heyting operations are preserved by inverse image in the first-order Heyting
fibration.  The modal steps are
Proposition~\ref{prop:dynamic-modal-beck-chevalley-span-comparison} and
Theorem~\ref{thm:dynamic-modal-invariance-trajectory-quotients}.  Tests and
guard restrictions are pullbacks, and quotient saturation ensures that their
comparison squares are cartesian.
\end{proof}

\begin{remark}[No automatic descent of arbitrary predicates]
\label{rem:dynamic-no-automatic-predicate-descent}
The result does not claim that every predicate on the original state object
descends through behavioural minimization.  Modal invariance is asserted for
predicates and guards already saturated with respect to the quotient.  This
qualification is necessary: a predicate that distinguishes states inside one
behavioural class cannot be represented on the quotient.
\end{remark}


\section{Specializations and examples}
\label{sec:dynamic-specializations-examples}

\subsection{One-object relative response logic}
\label{subsec:dynamic-one-object-response-logic}

Let \(M\) and \(N\) be monoid objects and let
\[
    \Phi:M\rightsquigarrow N
\]
have state object \(P\), right \(M\)-action \(x\cdot m\), and output cocycle
\[
    \omega_P:M\times P\to N.
\]
For a right \(N\)-object \(Y\), the relative state object is \(P\times Y\),
and the one-step transition is
\[
    (x,y)\xrightarrow{m}
    \left(
        x\cdot m,
        y\cdot\omega_P(m,x)
    \right).
\]
For predicates
\[
    \theta,\varphi\hookrightarrow P\times Y,
    \qquad
    R\hookrightarrow M,
    \qquad
    O\hookrightarrow N,
\]
the guarded diamond has the external reading
\[
\begin{aligned}
(x,y)\models\langle\theta;R/O\rangle\varphi
\quad\Longleftrightarrow\quad&
\theta(x,y)\ \text{and there exists }m\in M\text{ such that}\\
&
R(m),\quad O(\omega_P(m,x)),\\
&
\varphi\!\left(
x\cdot m,
y\cdot\omega_P(m,x)
\right).
\end{aligned}
\]
The corresponding strong profile box requires every \(R\)-input \(m\) to
produce an \(O\)-output and a continuation satisfying \(\varphi\).

\subsection{Stateless functors}
\label{subsec:dynamic-stateless-response-logic}

Let
\[
    F:\mathbb A\to\mathbb B
\]
be an internal functor, regarded as a stateless Mealy morphism.  The relative
state object is
\[
    A_0\times_{F_0,B_0,r}Y,
\]
and an admissible arrow \(a:u\to v\) produces
\[
    (v,y)\longmapsto
    \left(
        u,
        \sigma(F_1(a),y)
    \right).
\]
Thus every guarded modality is obtained from the general formulas by replacing
the state-dependent output \(\omega_P(a,x)\) with \(F_1(a)\).  In the terminal
downstream case the program category is
\(\mathbb A^{\mathrm{op}}\), and the output labelling is
\(F^{\mathrm{op}}\).

\subsection{Relational semantics in \(\mathbf{Set}\)}
\label{subsec:dynamic-relational-set-semantics}

A span
\[
    S\xleftarrow{\ell}R\xrightarrow{r}S
\]
determines the relation
\[
    x\,\overline R\,y
\]
when some \(\rho\in R\) satisfies
\[
    \ell(\rho)=x,
    \qquad
    r(\rho)=y.
\]
Then
\[
    x\models\langle R\rangle\varphi
\]
holds exactly when some \(\overline R\)-successor of \(x\) satisfies
\(\varphi\), and
\[
    x\models[R]\varphi
\]
holds exactly when every \(\overline R\)-successor of \(x\) satisfies
\(\varphi\).  In the relative Mealy case this relation lives on
\[
    S_{\Phi,Y}=P\times_{B_0}Y.
\]

At the profile level, the polynomial fibre
\[
    \mathcal M_{\mathbb A,\mathbb B,Y}(X)_{(v,y)}
    \cong
    \prod_{a:u\to v}
    \sum_{\beta:b'\to r(y)}
    X_{(u,\sigma(\beta,y))}
\]
is a total family of possible responses.  The existential lifting selects one
input position; the strong box constrains every guarded input position.

\subsection{Word inputs and generator-level paths}
\label{subsec:dynamic-word-generator-level-logic}

Let
\[
    \mathbb A=I^\ast
\]
be the one-object category of input words.  The response profile at a state is
total over all words:
\[
    w\longmapsto
    \left(
        \omega_P(w,x),
        \delta_P(w,x)
    \right).
\]
An unrestricted one-step modality may therefore quantify over an arbitrary
word, not only a letter.  Likewise, an unrestricted \(n\)-segment endpoint
modality collapses to the one-step modality by contraction.

Letter-sensitive logic is obtained in either of two ways:

\begin{enumerate}
\item restrict the one-step input guard to the length-one subobject
\(I\hookrightarrow I^\ast\);
\item begin with the generator-level execution graph and form its finite
trajectory objects before taking the free categorical completion.
\end{enumerate}

With either choice, a length-\(n\) trajectory records \(n\) letter steps and
does not collapse to a single letter transition.  This is the appropriate
setting for ordinary finite-word safety and reachability formulas.


\section{Chapter synthesis}
\label{sec:dynamic-response-trajectory-synthesis}

The three semantic presentations now fit into one logical comparison.

\begin{theorem}[Dynamic logic of relative responses and trajectories]
\label{thm:dynamic-response-trajectory-main-summary}
Let
\[
    \Phi:\mathbb A\rightsquigarrow\mathbb B
\]
be a Mealy morphism and \(Y\) a right \(\mathbb B\)-action.

\begin{enumerate}
\item The relative execution category supplies the one-step span
\[
    \mathbf m_{\Phi,Y}:
    S_{\Phi,Y}\nrightarrow S_{\Phi,Y}
\]
and hence span diamonds and, when available, span boxes.

\item Input, output, downstream, state, and interface predicates determine
guarded one-step subprograms
\[
    \theta;R/O/T.
\]

\item The response coalgebra supplies existential and universal profile
liftings.  Pullback along \(c_{\Phi,Y}\) identifies the existential lifting
with the guarded span diamond:
\[
    \theta\wedge
    c_{\Phi,Y}^\ast\lambda_{R/O/T}(\varphi)
    =
    \langle\theta;R/O/T\rangle\varphi.
\]

\item The weak profile box pulls back to the box of the restricted one-step
span:
\[
    c_{\Phi,Y}^\ast
    \Box^{\mathrm{prof}}_{R/O/T}(\varphi)
    =
    [R/O/T]\varphi.
\]
The strong profile box additionally requires every \(R\)-input position to
produce an \(O/T\)-guarded response.

\item Tests, sequential composition, finite guarded joins, conditionals,
domain predicates, and guarded Mealy automata form a generated span-program
calculus on \(S_{\Phi,Y}\).

\item Downstream action maps and evaluated Mealy simulations preserve the
modal semantics by Beck--Chevalley along their source-cartesian execution
squares.

\item The finite trajectory objects supply endpoint and boundary modalities.
Unrestricted endpoint horizons contract to one-step modalities, while
boundary-sensitive and generator-restricted formulas retain temporal
segmentation.

\item On the canonical behavioural machine, derivative generation is finite
endpoint reachability and the derivative safety core is universal
finite-path safety.

\item Strict regular quotients preserve the interpretation of predicates and
guards that descend to the quotient; in particular, quotient-saturated
formulas are invariant under behavioural minimization.
\end{enumerate}
\end{theorem}

\begin{proof}
Items (1)--(2) are the span and guarded one-step constructions of
Sections~\ref{sec:relative-predicate-transformers} and
\ref{sec:poly-relative-guard-loci-primitives}.  Items (3)--(4) are
Theorems~\ref{thm:poly-relative-lifting-span-agreement} and
\ref{thm:poly-relative-profile-box-span-agreement}.  Item (5) is developed in
Sections~\ref{sec:poly-relative-tests-choice}--
\ref{sec:poly-relative-generated-program-algebra}.  Item (6) is
Propositions~\ref{prop:relative-diamond-stability-downstream},
\ref{prop:relative-box-stability-downstream},
\ref{prop:poly-relative-diamond-stability-simulation}, and
\ref{prop:poly-relative-box-stability-simulation}.  Item (7) is
Propositions~\ref{prop:dynamic-finite-horizon-composition} and
\ref{prop:dynamic-unrestricted-endpoint-collapse}.  Item (8) combines
Theorem~\ref{thm:generated-closure-temporal-reachability} of the temporal
chapter with Theorem~\ref{thm:safety-core-universal-temporal-safety}.  Item
(9) is Theorem~\ref{thm:dynamic-modal-invariance-trajectory-quotients} and
Corollary~\ref{cor:dynamic-behavioural-invariance-descended-formulas}.
\end{proof}

\begin{remark}[Final conceptual summary]
\label{rem:dynamic-response-trajectory-final-summary}
The span, polynomial, and trajectory logics are not three unrelated modal
systems.  They inspect three presentations of the same relative Mealy
semantics:
\[
\begin{array}{ccl}
\text{total response profile}
&\xrightarrow{\ \text{flatten}\ }&
\text{one-step execution span}\\
&&\xrightarrow{\ \operatorname{Path}\ }
\text{finite trajectories}.
\end{array}
\]
Profile liftings express properties of a complete response policy.  Span
modalities express properties of selected executions.  Trajectory modalities
express properties of finite segmented executions.  Currying, flattening, and
the path construction explain exactly how these logical levels are related.
\end{remark}
